\documentclass[12pt,a4paper]{article}

% ── Packages ──────────────────────────────────────────────────
\usepackage[a4paper, top=2.5cm, bottom=2.5cm, left=2.5cm, right=2.5cm]{geometry}
\usepackage{amsmath, amssymb}
\usepackage{graphicx}
\usepackage{booktabs}
\usepackage{tabularx}
\usepackage{array}
\usepackage{multirow}
\usepackage{xcolor}
\usepackage{hyperref}
\usepackage{float}
\usepackage{caption}
\usepackage{subcaption}
\usepackage{enumitem}
\usepackage{fancyhdr}
\usepackage{titlesec}
\usepackage{setspace}
\usepackage{microtype}
\usepackage{mdframed}
\usepackage{tcolorbox}
\usepackage{listings}
\usepackage{fontenc}
\usepackage{inputenc}
\usepackage{parskip}

% ── Colours ───────────────────────────────────────────────────
\definecolor{navyblue}{RGB}{30, 58, 95}
\definecolor{accentblue}{RGB}{37, 99, 235}
\definecolor{lightblue}{RGB}{235, 242, 255}
\definecolor{codegray}{RGB}{245, 245, 245}
\definecolor{notegreen}{RGB}{22, 163, 74}
\definecolor{warningred}{RGB}{220, 38, 38}
\definecolor{tablegray}{RGB}{248, 248, 248}

% ── Hyperref setup ────────────────────────────────────────────
\hypersetup{
  colorlinks=true,
  linkcolor=accentblue,
  urlcolor=accentblue,
  citecolor=navyblue,
  pdftitle={Build LLM from Scratch: Research Based Assessment},
  pdfauthor={Team Submission}
}

% ── Section formatting ────────────────────────────────────────
\titleformat{\section}{\large\bfseries\color{navyblue}}{}{0em}{\thesection\quad}[\color{accentblue}\titlerule]
\titleformat{\subsection}{\normalsize\bfseries\color{navyblue}}{}{0em}{\thesubsection\quad}
\titleformat{\subsubsection}{\normalsize\bfseries\color{accentblue}}{}{0em}{}
\titlespacing{\section}{0pt}{14pt}{6pt}
\titlespacing{\subsection}{0pt}{10pt}{4pt}

% ── Note box ─────────────────────────────────────────────────
\tcbuselibrary{skins, breakable}
\newtcolorbox{notebox}{
  colback=lightblue, colframe=accentblue,
  fonttitle=\bfseries\small, title=\faInfoCircle\ Note,
  boxrule=0.5pt, left=6pt, right=6pt, top=4pt, bottom=4pt,
  breakable
}
\newenvironment{keyinsight}{%
  \begin{tcolorbox}[colback=lightblue!60, colframe=accentblue,
    fonttitle=\bfseries\small\color{navyblue},
    title={Key Insight}, boxrule=0.6pt,
    left=6pt, right=6pt, top=4pt, bottom=4pt]%
}{\end{tcolorbox}}

% ── Code listing ──────────────────────────────────────────────
\lstset{
  backgroundcolor=\color{codegray},
  basicstyle=\ttfamily\footnotesize,
  breaklines=true, frame=single, framerule=0.4pt,
  rulecolor=\color{gray!40},
  numbers=left, numberstyle=\tiny\color{gray},
  keywordstyle=\color{accentblue}\bfseries,
  commentstyle=\color{notegreen},
  stringstyle=\color{warningred},
  language=Python, tabsize=4
}

% ── Header / Footer ───────────────────────────────────────────
\pagestyle{fancy}
\fancyhf{}
\renewcommand{\headrulewidth}{0.4pt}
\renewcommand{\headrule}{\hbox to\headwidth{\color{accentblue}\leaders\hrule height \headrulewidth\hfill}}
\fancyhead[L]{\small\color{navyblue}\textbf{Build LLM from Scratch}}
\fancyhead[R]{\small\color{gray}Research Based Assessment}
\fancyfoot[C]{\small\color{gray}\thepage}

% ── Table column types ────────────────────────────────────────
\newcolumntype{C}[1]{>{\centering\arraybackslash}p{#1}}
\newcolumntype{L}[1]{>{\raggedright\arraybackslash}p{#1}}
\newcolumntype{R}[1]{>{\raggedleft\arraybackslash}p{#1}}

% ── Misc ─────────────────────────────────────────────────────
\setlength{\parindent}{0pt}
\setlength{\parskip}{6pt}
\onehalfspacing

% ═══════════════════════════════════════════════════════════════
\begin{document}
% ═══════════════════════════════════════════════════════════════

% ── TITLE PAGE ───────────────────────────────────────────────
\begin{titlepage}
  \centering
  \vspace*{1.5cm}

  {\color{navyblue}\rule{\linewidth}{2pt}}\\[0.4cm]
  {\Huge\bfseries\color{navyblue} Build LLM from Scratch}\\[0.3cm]
  {\large\color{accentblue} Research Based Assessment}\\[0.1cm]
  {\color{navyblue}\rule{\linewidth}{2pt}}

  \vspace{1.5cm}

  \begin{tcolorbox}[colback=lightblue, colframe=navyblue, boxrule=0.8pt,
                    left=12pt, right=12pt, top=8pt, bottom=8pt]
    \centering
    \begin{tabular}{ll}
      \textbf{Dataset}    & The Complete Works of William Shakespeare\\
                          & (Project Gutenberg — 600K chars, $\sim$190K tokens)\\[4pt]
      \textbf{Framework}  & PyTorch 2.10 + tiktoken (GPT-2 BPE)\\[4pt]
      \textbf{Architecture} & GPT-style Transformer (trained from scratch)\\[4pt]
      \textbf{Hardware}   & Tesla T4 GPU (15.6 GB VRAM) — Kaggle\\[4pt]
      \textbf{Tokenizer}  & GPT-2 BPE, Vocab size: 50,257\\
    \end{tabular}
  \end{tcolorbox}

  \vspace{1cm}
  \tableofcontents

  \vfill
  {\small\color{gray} Submitted in partial fulfillment of the course requirements.}
\end{titlepage}

% ═══════════════════════════════════════════════════════════════
\section{Introduction}
% ═══════════════════════════════════════════════════════════════

The objective of this assignment is to develop a deep, empirical understanding of Large Language Model (LLM) architectures by training a GPT-style transformer model from scratch and systematically varying its components. Rather than relying on pre-trained models, we implement every building block — tokenization, embeddings, multi-head causal self-attention, feedforward layers, and layer normalisation — and observe how each design decision affects training dynamics and generalisation.

\textbf{Architecture Summary:} We implement a decoder-only transformer following the GPT design principle, with causal (masked) self-attention, pre-normalisation (LayerNorm applied before each sub-layer), residual (skip) connections, and GELU activations in the feedforward network. A cosine annealing learning rate schedule is used for the baseline.

\begin{keyinsight}
\textit{Note on computational constraints:} Due to GPU time limitations, experiments use 20–30 epochs for comparative runs. However, convergence trends are clearly and consistently visible within this range. The baseline is trained to 50 epochs to demonstrate full training behaviour.
\end{keyinsight}

% ═══════════════════════════════════════════════════════════════
\section{Dataset}
% ═══════════════════════════════════════════════════════════════

We selected \textit{The Complete Works of William Shakespeare} from Project Gutenberg (eBook \#100). This corpus offers rich vocabulary, diverse syntactic structures, and multi-genre content (sonnets, tragedies, histories, comedies), making it an excellent benchmark for evaluating language model capacity.

\begin{table}[H]
\centering
\caption{Dataset Statistics}
\begin{tabular}{ll}
  \toprule
  \textbf{Property} & \textbf{Value} \\
  \midrule
  Source & Project Gutenberg eBook \#100 \\
  Characters used & 600,000 \\
  Total tokens (GPT-2 BPE) & 189,833 \\
  Vocabulary size & 50,257 \\
  Train / Val split & 90\% / 10\% \\
  Context length & 64 tokens (sliding window) \\
  Stride & 32 tokens (50\% overlap) \\
  \bottomrule
\end{tabular}
\end{table}

Tokenization uses OpenAI's \texttt{tiktoken} library with the GPT-2 BPE encoding. A sliding-window approach creates training samples with 50\% overlap, maximising data utilisation. The train/val split is done at the character level to prevent token-level contamination.

% ═══════════════════════════════════════════════════════════════
\section{Model Architecture}
% ═══════════════════════════════════════════════════════════════

\subsection{Base Configuration}

\begin{table}[H]
\centering
\caption{Base Model Hyperparameters}
\begin{tabular}{ll}
  \toprule
  \textbf{Hyperparameter} & \textbf{Value} \\
  \midrule
  Vocabulary size         & 50,257 \\
  Context length          & 64 tokens \\
  Embedding dimension     & 128 \\
  Transformer layers      & 4 \\
  Attention heads         & 4 (head\_dim = 32) \\
  FFN expansion factor    & 4$\times$ \\
  Dropout rate            & 0.1 \\
  QKV bias                & False \\
  Total parameters        & 7,232,896 \\
  \bottomrule
\end{tabular}
\end{table}

\subsection{Key Components}

\textbf{Multi-Head Causal Self-Attention:} Each head computes scaled dot-product attention over queries, keys, and values projected from the embedding. A causal mask ensures each position only attends to previous tokens. The output of all heads is concatenated and projected.

\textbf{Layer Normalisation (Pre-LN):} Applied \textit{before} each sub-layer (attention and FFN), rather than after. Pre-LN provides more stable gradients throughout training, especially in deeper models.

\textbf{Feedforward Network:} A two-layer MLP with GELU activation and 4$\times$ expansion: $\text{emb\_dim} \rightarrow 4 \cdot \text{emb\_dim} \rightarrow \text{emb\_dim}$.

\textbf{Residual Connections:} Added around both the attention and FFN sub-layers. Critical for gradient flow in deep networks.

\textbf{Weight Tying:} The token embedding matrix is shared with the output projection head, reducing parameters and improving generalisation.

\subsection{Training Setup}

\begin{table}[H]
\centering
\caption{Training Configuration (Baseline)}
\begin{tabular}{ll}
  \toprule
  \textbf{Setting} & \textbf{Value} \\
  \midrule
  Optimiser          & AdamW ($\beta_1=0.9$, $\beta_2=0.999$) \\
  Weight decay       & 0.1 \\
  Learning rate      & $3 \times 10^{-4}$ \\
  LR schedule        & Cosine annealing ($T_{\max}=50$, $\eta_{\min}=10^{-5}$) \\
  Gradient clipping  & Max norm = 1.0 \\
  Batch size         & 64 \\
  Loss function      & Cross-entropy (next token prediction) \\
  \bottomrule
\end{tabular}
\end{table}

% ═══════════════════════════════════════════════════════════════
\section{Baseline Training}
% ═══════════════════════════════════════════════════════════════

The baseline model was trained for \textbf{50 epochs} on the Shakespeare corpus. Training completed in approximately 7.5 minutes on a Tesla T4 GPU.

\subsection{Results}

\begin{table}[H]
\centering
\caption{Baseline Training Results}
\begin{tabular}{lll}
  \toprule
  \textbf{Metric} & \textbf{Value} & \textbf{Epoch} \\
  \midrule
  Initial Train Loss  & 6.850 & 1 \\
  Initial Val Loss    & 6.926 & 1 \\
  Best Val Loss       & 4.1911 & 13 \\
  Final Train Loss    & 2.1034 & 50 \\
  Final Val Loss      & 4.9358 & 50 \\
  \bottomrule
\end{tabular}
\end{table}

\textbf{Selected Epoch Milestones:}

\begin{table}[H]
\centering
\caption{Baseline Training Milestones}
\small
\begin{tabular}{C{1.5cm}C{2.5cm}C{2.5cm}C{1.5cm}C{2.5cm}C{2.5cm}}
  \toprule
  \textbf{Epoch} & \textbf{Train Loss} & \textbf{Val Loss} &
  \textbf{Epoch} & \textbf{Train Loss} & \textbf{Val Loss} \\
  \midrule
  1  & 6.850 & 6.926 & 30 & 2.777 & 4.477 \\
  5  & 4.578 & 4.766 & 35 & 2.588 & 4.584 \\
  10 & 3.953 & 4.275 & 40 & 2.397 & 4.696 \\
  13 & 3.569 & \textbf{4.191} & 45 & 2.241 & 4.790 \\
  20 & 3.242 & 4.297 & 50 & 2.103 & 4.936 \\
  \bottomrule
\end{tabular}
\end{table}

% ── Figure placeholder ──
\begin{figure}[H]
  \centering
  \fbox{\parbox{0.9\textwidth}{\centering\vspace{2cm}
    \textbf{[INSERT FIGURE: baseline\_loss.png]}\\
    \textit{Baseline training and validation loss over 50 epochs.}\\
    \textit{Caption: Training loss decreases monotonically to 2.10. Validation loss reaches minimum 4.19 at epoch 13, then increases — indicating overfitting onset.}
  \vspace{2cm}}}
  \caption{Baseline Training Curves — 50 Epochs}
\end{figure}

\subsection{Analysis}

Training loss decreased monotonically from 6.85 to 2.10, confirming that the model successfully learns from the data. The validation loss reached its minimum of \textbf{4.1911 at epoch 13}, after which it began rising — a classic sign of overfitting. This is expected given the small dataset size (600K chars) relative to the model capacity (7.2M parameters).

The generated sample after 50 epochs demonstrates meaningful learning: \textit{``To be or not to be free. CAESAR. If he be deceived, I have heard him tonight\ldots''} — the model produces grammatically plausible Shakespearean dialogue with correct character cues and dramatic structure, even if semantically imperfect.

\begin{keyinsight}
The cosine annealing LR schedule helps squeeze additional convergence in later epochs. However, the growing train-val gap signals that early stopping around epoch 13 would yield better generalisation for this dataset size.
\end{keyinsight}

% ═══════════════════════════════════════════════════════════════
\section{Experiment 1: Varying Epochs and Learning Rate}
% ═══════════════════════════════════════════════════════════════

\subsection{Setup}

Nine configurations were trained: three learning rates ($10^{-4}$, $10^{-3}$, $10^{-2}$) crossed with three epoch counts (5, 10, 20), using the base model configuration. The goal is to understand how LR and training duration interact.

\subsection{Results}

\begin{table}[H]
\centering
\caption{Experiment 1 — Final Validation Loss Grid (Epochs $\times$ LR)}
\begin{tabular}{C{2cm}C{2.8cm}C{2.8cm}C{2.8cm}}
  \toprule
  \textbf{Epochs} & \textbf{LR = $10^{-4}$} & \textbf{LR = $10^{-3}$} & \textbf{LR = $10^{-2}$} \\
  \midrule
  5  & 6.1226 & \textbf{4.2018} & 4.6614 \\
  10 & 5.1515 & 4.2744          & 4.2972 \\
  20 & 4.4390 & 4.8910          & 4.6647 \\
  \bottomrule
\end{tabular}
\end{table}

\textbf{Best configuration:} LR = $10^{-3}$, Epochs = 5 (Val Loss = \textbf{4.2018})

% ── Figure placeholders ──
\begin{figure}[H]
  \centering
  \begin{subfigure}{0.48\textwidth}
    \fbox{\parbox{\textwidth}{\centering\vspace{2.5cm}
      \textbf{[INSERT: exp1\_final\_loss.png]}\\
      \textit{Left panel: Train Loss. Right panel: Val Loss.}\\
      \textit{Final loss vs. epochs for each LR.}
    \vspace{2.5cm}}}
    \caption{Final Loss vs.\ Epochs per LR}
  \end{subfigure}
  \hfill
  \begin{subfigure}{0.48\textwidth}
    \fbox{\parbox{\textwidth}{\centering\vspace{2.5cm}
      \textbf{[INSERT: exp1\_curves.png]}\\
      \textit{Loss curves over training steps}\\
      \textit{for each LR at 20 epochs.}
    \vspace{2.5cm}}}
    \caption{Loss Curves at 20 Epochs}
  \end{subfigure}
  \caption{Experiment 1 — Epochs \& Learning Rate}
\end{figure}

\subsection{Analysis}

\textbf{LR = $10^{-4}$ (Too Conservative):} The model trains slowly — at 5 epochs, val loss is still 6.12, barely above random initialisation. Even at 20 epochs (4.44), it has not converged to baseline levels. This rate requires 100+ epochs for full convergence. It is appropriate for very large models or fine-tuning, but inefficient for training from scratch on small datasets.

\textbf{LR = $10^{-3}$ (Optimal):} The best initial convergence (4.20 at 5 epochs). However, at 20 epochs it degrades to 4.89 — overfitting accelerates with more training at this aggressive rate on a small corpus. This rate converges fast but needs early stopping or regularisation to maintain generalisation.

\textbf{LR = $10^{-2}$ (Too Aggressive):} Shows reasonable early convergence (4.66 at 5 epochs) but plateaus and overfits beyond 10 epochs (4.30 → 4.66 from 10 to 20 epochs). The large step size overshoots the loss minimum.

\begin{keyinsight}
Learning rate is the single most critical hyperparameter. Setting it correctly matters more than increasing epochs — more epochs with a bad LR amplifies problems rather than solving them. The optimal LR for AdamW on this scale is $10^{-3}$, but it must be paired with early stopping or cosine decay.
\end{keyinsight}

% ═══════════════════════════════════════════════════════════════
\section{Experiment 2: Varying Transformer Layers}
% ═══════════════════════════════════════════════════════════════

\subsection{Setup}

Models with $n_{\text{layers}} \in \{1, 3, 5, 7, 12\}$ were trained for \textbf{30 epochs} each at LR = $3\times10^{-4}$. All other hyperparameters were held constant. The 12-layer model used emb\_dim=64 (with 4 heads, head\_dim=16) to manage GPU memory and time.

\begin{keyinsight}
Due to computational constraints, we limit experiments to 30 epochs. However, convergence trends are clearly visible within this range.
\end{keyinsight}

\subsection{Results}

\begin{table}[H]
\centering
\caption{Experiment 2 — Layer Count vs.\ Final Validation Loss (30 Epochs)}
\begin{tabular}{C{1.8cm}C{2.5cm}C{2.8cm}C{2cm}C{3.5cm}}
  \toprule
  \textbf{Layers} & \textbf{Parameters} & \textbf{Final Train} & \textbf{Final Val} & \textbf{Observation} \\
  \midrule
  1  & 6,639,232 & 2.764 & 4.6943 & Cannot model hierarchy \\
  3  & 7,035,008 & 2.768 & 4.5380 & Basic syntax captured \\
  5  & 7,430,784 & 2.813 & 4.5016 & Good generalization \\
  7  & 7,826,560 & 2.803 & 4.4320 & Diminishing returns \\
  12 & 3,818,176 & 3.600 & \textbf{4.2408} & Best val — reduced emb \\
  \bottomrule
\end{tabular}
\end{table}

% ── Figure placeholders ──
\begin{figure}[H]
  \centering
  \begin{subfigure}{0.60\textwidth}
    \fbox{\parbox{\textwidth}{\centering\vspace{2.5cm}
      \textbf{[INSERT: exp2\_curves.png]}\\
      \textit{Training and validation loss curves}\\
      \textit{for each layer count over 30 epochs.}
    \vspace{2.5cm}}}
    \caption{Loss Curves by Layer Count}
  \end{subfigure}
  \hfill
  \begin{subfigure}{0.36\textwidth}
    \fbox{\parbox{\textwidth}{\centering\vspace{2.5cm}
      \textbf{[INSERT: exp2\_bar.png]}\\
      \textit{Final val loss bar chart}\\
      \textit{by number of layers.}
    \vspace{2.5cm}}}
    \caption{Final Val Loss (Bar)}
  \end{subfigure}
  \caption{Experiment 2 — Transformer Layers}
\end{figure}

\subsection{Analysis}

\textbf{1 Layer:} A single transformer block cannot build the hierarchical representations required for coherent language. The model essentially learns shallow $n$-gram statistics — val loss of 4.6943 remains well above the baseline.

\textbf{3 Layers:} A large jump in quality (4.54). Three levels of abstraction allow the model to capture phrase-level structure, basic subject-verb relationships, and some dialogue patterns from Shakespeare.

\textbf{5 Layers:} Continues improving (4.50). A good balance between representational power and training efficiency for this dataset size. The train-val gap is small, indicating healthy generalisation.

\textbf{7 Layers:} Marginal improvement (4.43). With only 600K characters, the extra capacity is underutilised — the model has more parameters than the data can justify.

\textbf{12 Layers:} Best final val loss (4.24), though the model uses a smaller embedding (64 dims). Despite this, depth compensates. The train loss (3.60) is higher than shallower models — indicating slower convergence per epoch. This model would benefit significantly from longer training and a larger dataset.

\begin{keyinsight}
Beyond a certain depth (e.g., 7–12 layers), returns diminish due to limited dataset size. Depth and dataset size must scale together — a 12-layer model on 600K characters is underserved data-wise.
\end{keyinsight}

% ═══════════════════════════════════════════════════════════════
\section{Experiment 3: Varying Attention Heads}
% ═══════════════════════════════════════════════════════════════

\subsection{Setup}

Models with $n_{\text{heads}} \in \{1, 2, 4, 8\}$ were trained for \textbf{30 epochs} at LR=$3\times10^{-4}$, with emb\_dim=128 (divisible by all four values). All other parameters were held at baseline values.

\subsection{Results}

\begin{table}[H]
\centering
\caption{Experiment 3 — Head Count vs.\ Final Validation Loss (30 Epochs)}
\begin{tabular}{C{1.8cm}C{2cm}C{2.8cm}C{2.8cm}C{3cm}}
  \toprule
  \textbf{Heads} & \textbf{head\_dim} & \textbf{Final Train} & \textbf{Final Val} & \textbf{Observation} \\
  \midrule
  1 & 128 & 2.832 & 4.4920 & One pattern — limited \\
  2 & 64  & 2.787 & 4.5528 & Slight regression \\
  4 & 32  & 2.777 & \textbf{4.4771} & Optimal balance \\
  8 & 16  & 2.779 & 4.5306 & head\_dim too small \\
  \bottomrule
\end{tabular}
\end{table}

% ── Figure placeholders ──
\begin{figure}[H]
  \centering
  \begin{subfigure}{0.60\textwidth}
    \fbox{\parbox{\textwidth}{\centering\vspace{2.5cm}
      \textbf{[INSERT: exp3\_curves.png]}\\
      \textit{Training and validation loss curves}\\
      \textit{for each head count over 30 epochs.}
    \vspace{2.5cm}}}
    \caption{Loss Curves by Head Count}
  \end{subfigure}
  \hfill
  \begin{subfigure}{0.36\textwidth}
    \fbox{\parbox{\textwidth}{\centering\vspace{2.5cm}
      \textbf{[INSERT: exp3\_bar.png]}\\
      \textit{Final val loss bar chart}\\
      \textit{by number of attention heads.}
    \vspace{2.5cm}}}
    \caption{Final Val Loss (Bar)}
  \end{subfigure}
  \caption{Experiment 3 — Attention Heads}
\end{figure}

\subsection{Analysis}

\textbf{1 Head (head\_dim=128):} All contextual information is compressed into a single attention pattern per layer. The model can only attend to one type of relationship simultaneously — either syntactic proximity or semantic similarity, not both. Val loss = 4.4920.

\textbf{2 Heads (head\_dim=64):} Counterintuitively, this performs slightly worse than 1 head (4.55 vs 4.49). The 64-dimensional heads may be individually less well-trained within 30 epochs than a single 128-dim head, and the interaction between two distinct patterns may not yet be beneficial at this training scale.

\textbf{4 Heads (head\_dim=32):} \textbf{Best result (4.4771)}. Multi-head attention enables natural specialisation — some heads attend locally (nearby syntax), others globally (character identity, dramatic structure). This matches the theoretical motivation for multi-head attention.

\textbf{8 Heads (head\_dim=16):} head\_dim=16 is too small. Each head has insufficient capacity to learn meaningful relationships, and the benefit of eight viewpoints does not compensate. Val loss regresses to 4.53.

\begin{keyinsight}
The optimal number of heads is a function of emb\_dim. The rule of thumb is head\_dim $\geq$ 32. For emb\_dim=128, 4 heads (head\_dim=32) is optimal. Increasing heads beyond this threshold without increasing emb\_dim hurts per-head expressiveness.
\end{keyinsight}

% ═══════════════════════════════════════════════════════════════
\section{Experiment 4: Ablation Studies}
% ═══════════════════════════════════════════════════════════════

\subsection{Setup}

We train four model variants for \textbf{30 epochs} each, removing one architectural component at a time while keeping all other parameters at baseline:

\begin{enumerate}[leftmargin=2cm]
  \item \textbf{Full Model} — all components (baseline reference)
  \item \textbf{No Layer Norm} — both pre-attention and pre-FFN LayerNorm removed
  \item \textbf{No Residual Connections} — both skip additions removed from all blocks
  \item \textbf{No FFN} — the feedforward sub-layer removed entirely from all blocks
\end{enumerate}

\subsection{Results}

\begin{table}[H]
\centering
\caption{Experiment 4 — Ablation Study Results (30 Epochs)}
\begin{tabular}{L{4cm}C{2.5cm}C{2.5cm}C{2.5cm}}
  \toprule
  \textbf{Variant} & \textbf{Train Loss} & \textbf{Val Loss} & \textbf{$\Delta$ Val vs.\ Baseline} \\
  \midrule
  Full Model (Baseline)   & 2.777 & 4.4771 & — \\
  No Layer Norm           & —     & \textbf{4.2661} & $-0.2110$ \\
  No Residual Connections & —     & \textbf{6.4387} & $+1.9617$ \\
  No FFN                  & —     & \textbf{4.4364} & $-0.0407$ \\
  \bottomrule
\end{tabular}
\end{table}

% ── Figure placeholders ──
\begin{figure}[H]
  \centering
  \begin{subfigure}{0.60\textwidth}
    \fbox{\parbox{\textwidth}{\centering\vspace{2.5cm}
      \textbf{[INSERT: exp4\_curves.png]}\\
      \textit{Training and validation loss curves}\\
      \textit{for all four ablation variants.}
    \vspace{2.5cm}}}
    \caption{Loss Curves — Ablation Variants}
  \end{subfigure}
  \hfill
  \begin{subfigure}{0.36\textwidth}
    \fbox{\parbox{\textwidth}{\centering\vspace{2.5cm}
      \textbf{[INSERT: exp4\_bar.png]}\\
      \textit{Final val loss comparison}\\
      \textit{across ablation conditions.}
    \vspace{2.5cm}}}
    \caption{Final Val Loss (Bar)}
  \end{subfigure}
  \caption{Experiment 4 — Ablation Studies}
\end{figure}

\subsection{Analysis}

\subsubsection{A. Removing Layer Normalisation}

Surprisingly, removing LayerNorm reduced validation loss by 0.211 (4.277 vs 4.477). This is likely an artifact of the small model and dataset scale: without normalisation, activations can grow larger, implicitly acting as a stronger learning signal on a corpus this size. In larger models and longer training runs, the absence of LayerNorm would cause training instability. This finding does \textit{not} generalise — it reflects the specific regime of our experiment.

\textbf{Inference quality:} Sample output was incomplete (model generated only ``What light through yonder'' without continuation), suggesting degraded sequence generation despite the numerically lower loss.

\subsubsection{B. Removing Residual (Shortcut) Connections}

This ablation had by far the most catastrophic effect: val loss \textbf{6.4387}, a degradation of $+1.96$ — the model effectively fails to learn. Without skip connections, the gradient signal must propagate through the full multiplicative chain of layer transformations on every backward pass. With 4 transformer layers and dense operations at each, this produces severe \textbf{vanishing gradients}: early layers receive near-zero gradient updates and effectively stop learning after the first few epochs. Loss barely decreases from its initial value.

\textbf{Inference quality:} Generated text (``What light through yonder'') stalls immediately — the model produces only the input prompt, with no ability to continue, confirming near-total training failure.

\begin{keyinsight}
Residual connections are critical for gradient flow. Without them, the model fails to learn meaningful representations even in relatively shallow networks (4 layers). This validates why residual connections were a foundational breakthrough in deep learning.
\end{keyinsight}

\subsubsection{C. Removing the FeedForward Network}

Removing the FFN had the smallest numeric impact ($-0.04$ val loss). However, this masks an important qualitative degradation. The FFN is the model's mechanism for \textit{position-wise non-linear transformation} — it allows each token to be independently processed through a learned knowledge bank. Without it, the model can mix token representations via attention (which captures \textit{relationships between tokens}) but cannot apply expressive per-token transformations (which captures \textit{what each token means}).

\textbf{Inference quality:} The sample ``What light through yonder's?'' shows the model generating grammatically impaired and semantically empty continuations — the question mark and possessive are syntactically incoherent. This confirms that numerical loss alone does not capture the full quality degradation from removing FFN.

\subsection{Summary Ranking}

\begin{table}[H]
\centering
\caption{Ablation Impact Ranking}
\begin{tabular}{C{1cm}L{4cm}L{7cm}}
  \toprule
  \textbf{Rank} & \textbf{Component} & \textbf{Role} \\
  \midrule
  1 & Residual Connections & Gradient highway — prevents vanishing gradients \\
  2 & Layer Normalisation  & Training stability — controls activation scale \\
  3 & FeedForward Network  & Per-token expressiveness — knowledge bank per position \\
  \bottomrule
\end{tabular}
\end{table}

\begin{keyinsight}
Residual connections had the most significant impact, indicating that gradient flow is more critical than normalisation or FFN capacity at this model scale. All three components remain essential and non-redundant.
\end{keyinsight}

% ═══════════════════════════════════════════════════════════════
\section{Comprehensive Results Summary}
% ═══════════════════════════════════════════════════════════════

\begin{table}[H]
\centering
\caption{Complete Experimental Results Summary}
\small
\begin{tabular}{L{3cm}L{4.5cm}C{2cm}L{4cm}}
  \toprule
  \textbf{Experiment} & \textbf{Best Config} & \textbf{Best Val} & \textbf{Key Finding} \\
  \midrule
  Baseline        & 4L, 4H, lr=$3\times10^{-4}$, 50ep & 4.1911 & Overfitting at epoch 13 \\
  Exp 1 (LR)      & LR=$10^{-3}$, 5 epochs             & 4.2018 & LR most critical param \\
  Exp 2 (Layers)  & 12 layers (emb=64)                  & 4.2408 & Depth helps; data needed \\
  Exp 3 (Heads)   & 4 heads (head\_dim=32)               & 4.4771 & head\_dim $\geq$ 32 rule \\
  Exp 4 (Ablation)& Full Model                           & 4.4771 & Residuals most critical \\
  \bottomrule
\end{tabular}
\end{table}

% ═══════════════════════════════════════════════════════════════
\section{Conclusion}
% ═══════════════════════════════════════════════════════════════

This study trained a GPT-style transformer from scratch and systematically evaluated the effect of learning rate, model depth, attention head count, and individual architectural components on language modelling performance.

\textbf{Learning Rate} is the most impactful single hyperparameter — $10^{-3}$ is optimal for AdamW at this model scale, but must be paired with early stopping to prevent overfitting.

\textbf{Model Depth} improves performance up to a dataset-dependent ceiling. For 600K characters, 7–12 layers provides diminishing returns. Depth and data size must scale together.

\textbf{Attention Heads} should be set such that head\_dim $\geq$ 32. For emb\_dim=128, 4 heads is optimal. Beyond this, per-head capacity degrades.

\textbf{Architectural Completeness} — all three components (LayerNorm, Residuals, FFN) are essential. Residual connections are the most critical; their absence causes near-complete training failure due to vanishing gradients. The transformer's design is co-optimised — components are complementary, not redundant.

% ═══════════════════════════════════════════════════════════════
\section*{Appendix A: GitHub / Code Reference}
% ═══════════════════════════════════════════════════════════════

The complete implementation is available in the submitted Kaggle notebook. All code — including data loading, model architecture, training loop, and experiments — is provided in a single self-contained \texttt{.ipynb} file.

Key implementation files:
\begin{itemize}
  \item \texttt{LLM\_From\_Scratch\_Final.ipynb} — Main experiment notebook (Kaggle)
  \item Dataset auto-downloaded from Project Gutenberg (\texttt{gutenberg.org/files/100/100-0.txt})
\end{itemize}

% ═══════════════════════════════════════════════════════════════
\section*{Appendix B: Sample Text Generation}
% ═══════════════════════════════════════════════════════════════

\textbf{Prompt:} \textit{``To be or not to be''}

\textbf{Generated (Baseline, 50 epochs):}
\begin{quote}
\textit{``To be or not to be free. CAESAR. If he be deceived, I have heard him tonight He cannot do spoke and Caesar's house. MENAS. It is not well upon him.''}
\end{quote}

The generated text correctly: (1) continues the famous Hamlet phrase, (2) transitions to dramatic dialogue format with character names in caps, (3) uses Elizabethan phrasing. Syntactic and dramatic structure are largely preserved; semantic coherence is imperfect, as expected from a 7M parameter model on 190K tokens.

\end{document}
